% Corrected French for printed p.163.
% Corrections: flasquee -> flasque; suite exact -> suite exacte;
% si ils -> s'ils.
\SourcePage{163}

\LayerTarget{tohoku:claim:3.4.2:H1-bijection:surjectivity}
si $L$ est localement isomorphe à $\overline F$ (car on voit facilement
qu'un faisceau localement isomorphe à un faisceau flasque est flasque).
C'est un épimorphisme, car soit $(h_{ij})$ un 1-cocycle de $H$ relatif à un
recouvrement ouvert $(U_i)$ de $X$, comme $H^1(X,\overline H)$ est réduit à
l'élément neutre, il existe des $\overline h_i\in\Gamma(U_i,\overline H)$
telles que $h_{ij}=\overline h_i^{-1}\overline h_j$, d'autre part on peut
évidemment relever les $\overline h_i$ en des sections $\overline g_i$ de
$\overline G$ (il suffit de se rappeler de la définition de $\overline G$ et
$\overline H$ !). Posant $p_{ij}=\overline g_i^{-1}\overline g_j$, on voit
aussitôt que la section de $H$ sur $U_{ij}$ définie par $p_{ij}$ est $h_{ij}$,
d'où on conclut que $p_{ij}\in\Gamma(U_{ij},P)$, donc $(p_{ij})$ est un
1-cocycle de $P$ définissant le 1-cocycle donné par passage au quotient.

\LayerTarget{tohoku:construction:3.4.2:associated-P-bundle}
Du résultat précédent on conclut que le fibré $E$ est isomorphe au fibré
associé à un fibré $Q$ de faisceau structural $P$, bien déterminé d'ailleurs
à isomorphisme près.

\LayerTarget{tohoku:action:3.4.2:P-actions}
Notons d'ailleurs que faisant opérer
$P=\overline F G\subset\overline G$ sur $\overline G$ par automorphismes
intérieurs, $F$ reste stable par ces opérations (puisque $\overline F$ opère
trivialement sur $\overline F$, et que $F$ est invariant dans $G$), donc $P$
opère dans $\overline F$ de façon naturelle. De plus $F$ est stable sous les
opérations de $P$, et les opérations de $P$ sur $F$ ainsi obtenues ne sont
autres que celles qu'on obtient en composant $P\longrightarrow H$ et la
représentation naturelle de $H$ par des opérations dans $F$. De ceci on
déduit donc que le faisceau associé $F^E$ s'identifie aussi au faisceau
associé $F^Q$. Par ailleurs, les opérations de $P$ passent aussi au quotient
$F'=\overline F/F$ ; les représentations de $P$ par germes d'automorphismes
de $F,\overline F,F'$ seront notées par $\sigma$. Notons maintenant qu'on a
une suite exacte d'homomorphismes de faisceaux :

\LayerTarget{tohoku:eq:3.4.2:crossed-extension}
\[
e\longrightarrow G\longrightarrow P\mathop{\longrightarrow}^{u}F'
\longrightarrow e.
\]

\LayerTarget{tohoku:map:3.4.2:u}
L'homomorphisme $G\longrightarrow P=\overline F G$ est l'homomorphisme
d'injection, donc un homomorphisme de faisceaux de groupes, et
l'homomorphisme $u:P\longrightarrow F'$ est défini en faisant correspondre,
à un produit $\overline f g$ ($\overline f\in\overline F(x)$,
$g\in G(x)$) la classe de $\overline f$ dans $F'(x)$ (cette définition a un
sens, grâce au fait que $G\cap\overline F=F$).

\LayerTarget{tohoku:claim:3.4.2:u-properties}
Cet homomorphisme \textit{ne respecte pas} en général les structures
multiplicatives, mais satisfait aux conditions suivantes : (i) $u$ est
surjectif, et deux éléments de $P$ ont la même image si et seulement s'ils
sont congrus sous $G$ opérant à droite (i. e. s'ils définissent le même
élément de $P/G$) (ii) $u$ est un homomorphisme croisé de $P$ dans $F'$
($P$ opérant sur $F'$ comme indiqué plus haut), i. e. on a, dans chaque fibre
$P(x)$, $u(e)=e$ et $u(pp')=u(p)\sigma(p)u(p')$.

\LayerTarget{tohoku:construction:3.4.2:dQ-obstruction}
Nous allons de cette situation, et de la donnée du fibré $Q$ à faisceau
structural $P$, déduire un élément $d(Q)\in H^1(X,{F'}^Q)$ (où ${F'}^Q$ est
le faisceau associé à $Q$ et aux opérations données de $P$ sur $F'$), de telle
façon que la nullité de $d(Q)$ soit nécessaire et suffisante pour que la
classe $c(Q)\in H^1(X,P)$ de $Q$ soit dans l'image de $H^1(X,G)$. Alors la
proposition 3.4.2 sera prouvée (modulo des vérifications immédiates de
naturalité).

\LayerTarget{tohoku:reduction:3.4.2:E-via-Q}
En effet la condition trouvée est nécessaire et suffisante pour que la classe
$c(E)\in H^1(X,H)$ soit dans l'image de $H^1(X,G)$ (comme on voit aussitôt
sur le diagramme commutatif

\TohokuContinuation{164}
